\chapter{Primary Angle-Closure Glaucoma}
\index{Primary Angle-Closure Glaucoma}
\label{sec:Primary Angle-Closure Glaucoma}

\section{Overview}


\begin{itemize}[noitemsep]
  \item Primary angle-closure glaucoma is an eye emergency that occurs when the peripheral iris physically obstructs the trabecular meshwork, blocking aqueous humor removal of fluid and causing a rapid rise in intraocular pressure
  \item It is more common in hyperopic eyes with shallow anterior chambers.
\end{itemize}
\section{Causes}
This condition happens because of anatomical predisposition leading to iridotrabecular contact.   Degenerative processes involve progressive cellular dysfunction and tissue breakdown over time. Contributing factors include oxidative stress, inflammation, accumulated cellular damage, and reduced regenerative capacity.  Protein aggregation and accumulation is a common mechanism in many neurodegenerative disorders.  Genetic mutations account for some neurological conditions, while others are sporadic or environmentally triggered.   
\section{Risk Factors}
\begin{itemize}[noitemsep]
  \item Genetic predisposition and family history
  \item Advancing age
  \item Lifestyle factors (diet, exercise, smoking, alcohol)
  \item Environmental exposures
  \item Underlying medical conditions
  \item Medication use
\end{itemize}
\section{Signs and Symptoms}
\subsection{Common symptoms}
\begin{itemize}[noitemsep]
  \item \item Headache and facial pain Weakness or paralysis of muscles Numbness, tingling, or abnormal sensations Vision changes including blurred or double vision Difficulty with balance and coordination Cognitive changes (memory loss, confusion, difficulty concentrating) Speech and language difficulties Seizures or involuntary movements Dizziness and vertigo Fatigue and sleep disturbances
  \item Pain or discomfort in the affected area
  \end{itemize}
\subsection{Emergency warning signs}
\begin{itemize}[noitemsep]
  \item Severe or worsening symptoms of primary angle-closure glaucoma require immediate medical evaluation
\end{itemize}
\section{Progression and Stages}
The condition may progress through different stages of severity over time. Early stages often involve mild or subtle symptoms that may be overlooked. Moderate stages involve more noticeable symptoms affecting daily function. Advanced stages may involve significant impairment and complications.

\begin{itemize}[noitemsep]
  \item Symptoms of acute angle closure include severe eye pain, headache, nausea, blurred vision, halos around lights, a mid-dilated fixed pupil, and markedly elevated IOP
  \item Chronic angle closure may be asymptomatic until advanced.
\end{itemize}
\section{Complications}
\begin{itemize}[noitemsep]
  \item Disease progression leading to worsening symptoms
  \item Secondary infections or injuries
  \item Side effects of treatment
  \item Reduced quality of life
  \item Emotional and psychological impact
\end{itemize}
\section{Diagnosis}
\begin{itemize}[noitemsep]
  \item Doctors diagnose this based on tonometry and gonioscopy. Treatment for acute angle closure includes immediate topical and systemic IOP-lowering agents followed by laser peripheral iridotomy
  \item Fellow eyes should undergo prophylactic iridotomy.
  \item Comprehensive medical history and physical examination
  \item Laboratory tests to assess organ function
  \item Imaging studies to visualize affected structures
  \item Specialized testing depending on the affected system
  \item Consultation with relevant specialists
\end{itemize}
\section{Prevention}
\begin{itemize}[noitemsep]
  \item Maintain a healthy lifestyle with balanced diet and exercise
  \item Avoid known risk factors
  \item Attend regular medical check-ups for early detection
  \item Follow recommended screening guidelines
\end{itemize}
\section{Treatment Options}
Pharmacologic therapy targeting specific neurotransmitter systems or disease mechanisms Physical, occupational, and speech therapy for functional maintenance Surgical interventions when appropriate (deep brain stimulation, tumor resection) Cognitive rehabilitation and memory support Pain management for neuropathic pain Assistive devices and mobility aids Lifestyle modifications including exercise, nutrition, and sleep hygiene Support services and caregiver education Regular neurologic monitoring and treatment adjustment Clinical trials for emerging therapies

\begin{itemize}[noitemsep]
  \item Medications to manage symptoms and address underlying mechanisms
  \item Lifestyle modifications and self-care measures
  \item Physical therapy and rehabilitation
  \item Surgical intervention when indicated
  \item Regular monitoring and follow-up care
\end{itemize}
\section{Living with It}
\begin{itemize}[noitemsep]
  \item Follow your treatment plan as prescribed
  \item Maintain regular follow-up with your healthcare provider
  \item Make necessary lifestyle adjustments
  \item Seek support from family, friends, or support groups
  \item Stay informed about your condition
\end{itemize}
\section{Outlook}
Most people recover well with prompt treatment. Neurological conditions vary widely in their course and outcomes. Some are progressive, while others follow a relapsing-remitting pattern. Early intervention and rehabilitation can improve outcomes. Neurological conditions vary significantly in their course, from acute and self-limited to chronic and progressive. Early diagnosis and intervention can slow progression and improve quality of life. Rehabilitation and supportive therapies play crucial roles in maintaining function. Research continues to develop disease-modifying treatments for previously untreatable conditions. Multidisciplinary care addressing medical, functional, and psychosocial needs optimizes outcomes.
